\documentclass[12pt]{article}
\usepackage[utf8]{inputenc}
\usepackage[T1]{fontenc}
\usepackage{hyperref}
\title{Kua-nsi Language Resources for K-12 Curriculum Development}
\author{Compiled by Claude Sonnet 5.5 (claude-sonnet-5-5) for Andr\'as Kornai}
\date{}
\begin{document}
\maketitle

\section{Introduction}
Kua-nsi (ISO 639-3 ykn, Glottocode kuan1249) is a Loloish language of the Central Ngwi (Yi) group of Sino-Tibetan. It is spoken by roughly 5,000 to 5,400 people in Liuhe Township, Heqing County, Dali Bai Autonomous Prefecture, Yunnan, China. Speakers call themselves Kua'eshi, and their Han and Bai neighbours call them Baiyi. Children still learn it as a first language, but Mandarin is increasingly dominant among the young and schooling is in Standard Mandarin. The search found no Kua-nsi school materials.

\section{Primary Educational Resources}
The searches found no primers, textbooks or school programmes in Kua-nsi.

\section{Linguistic and Reference Materials}
\begin{itemize}
\item Wikipedia, ``Kua-nsi language'': overview of classification and phonology. \url{https://en.wikipedia.org/wiki/Kua-nsi_language}
\item ``IPA illustration of Kua-nsi'' (Journal of the International Phonetic Association, Cambridge Core, PDF): a phonetic description of the language. \url{https://core-cms.cambridgecore.org/core/services/aop-cambridge-core/content/view/618FA2E06E6EA639FDED569F86718003/S0025100324000161a.pdf/ipa-illustration-of-kuansi.pdf}
\item Max Planck Institute (EVA) Numerals in Languages of the Yi group: Kua-nsi Yi, Yunnan. \url{https://lingweb.eva.mpg.de/channumerals/Yi-Kua-nsi.htm}
\end{itemize}

\section{Religious and Community Resources}
\begin{itemize}
\item Joshua Project, Kua-nsi people-group profile. \url{https://joshuaproject.net/people_groups/20588}
\item Joshua Project, language page for ykn. \url{https://joshuaproject.net/index.php/languages/ykn}
\end{itemize}
I found no community organisations and no scripture in the language.

\section{Digital and Audio Resources}
The searches found no audio, video or corpus resources.

\section{Curriculum Development Context}
The language is stable and used by all community members, which is a good base. Instruction is in Mandarin, however, and there is no known Kua-nsi orthography in school use. The IPA illustration and similar linguistic descriptions could support the design of a practical orthography. Further materials would need to be created with the community.

\end{document}